% Source PDF: source/普通高考/2012/2012广东文.pdf
% Page: 1 (question 17)
% PDF embedded bitmap attempt: pdfimages -list found no image objects; the source page is vector text/line art.
% Grid evidence: tmp/repaint_2012_guangdong_liberal/source_pages/page-1-grid.png and q17_grid_crop3.png.
% Original comparison crop: tmp/repaint_2012_guangdong_liberal/q17_orig_crop.png, cropped from the no-grid source-page render.
% Elements: score axis with origin 0, an axis break before 50, ticks 50--100, frequency-density axis,
% five contiguous bins [50,60), [60,70), [70,80), [80,90), [90,100], and dashed guides at a, 0.02, 0.03, 0.04.
% Relations: the bar heights are (a, 0.04, 0.03, 0.02, a); since all class widths are 10 and the 100 frequencies sum to 1,
% 10(2a+0.04+0.03+0.02)=1, hence a=0.005.
% solid segments: bar outlines, x/y axis arrows, tick marks, and the x-axis break stroke.
% dashed segments: horizontal reference lines at the four source-labelled heights.
% labels: 频率/组距 sits above the y-axis, 成绩 follows the x-axis arrow, and source tick labels stay below/left of axes.
\begin{tikzpicture}[baseline,font=\normalsize]
\begin{axis}[exam statistical histogram,
    width=10.4cm,
    height=5.8cm,
    axis lines=none,
    clip=false,
    % Internal score coordinates are equally spaced: 0 is the origin and
    % 1--6 are the displayed group boundaries 50--100.  This preserves one
    % group width for the broken 0-to-50 segment.
    xmin=-2.0,
    xmax=8.0,
    ymin=-0.005,
    ymax=0.053,
    ticks=none,
]
    % Source dashed guides terminate at the corresponding histogram boundary.
    \draw[dashed,line width=0.55pt] (axis cs:0,0.005) -- (axis cs:6,0.005);
    \draw[dashed,line width=0.55pt] (axis cs:0,0.02) -- (axis cs:4,0.02);
    \draw[dashed,line width=0.55pt] (axis cs:0,0.03) -- (axis cs:3,0.03);
    \draw[dashed,line width=0.55pt] (axis cs:0,0.04) -- (axis cs:2,0.04);

    % Frequency-density bars read from the source histogram.
    \addplot[
        ybar interval,
        draw=black,
        fill=white,
        line width=0.7pt,
    ] coordinates {
        (1,0.005)
        (2,0.04)
        (3,0.03)
        (4,0.02)
        (5,0.005)
        (6,0)
    };

    % Source x/y axes and the short break before the score bins.
    \draw[->,line width=0.75pt] (axis cs:0,0) -- (axis cs:7,0);
    \draw[->,line width=0.75pt] (axis cs:0,0) -- (axis cs:0,0.050);
    \examstatxbreak[line width=0.75pt]{0.24000000000000002}{0};

    \examstatylabel{\(\dfrac{\text{频率}}{\text{组距}}\)};
    \examstatxlabel{成绩};

    % Tick marks and labels reproduced from the source chart.
    \draw[line width=0.45pt] (axis cs:0,0) -- (axis cs:0,-0.0010);
    \node[anchor=north east,inner sep=0pt,font=\normalsize,yshift=-4pt] at (axis cs:-0.25,0.000000) {0};
    \draw[line width=0.45pt] (axis cs:1,0) -- (axis cs:1,-0.0010);
    \node[anchor=north,inner sep=0pt,font=\normalsize,yshift=-2pt] at (axis cs:1,-0.001000) {50};
    \draw[line width=0.45pt] (axis cs:2,0) -- (axis cs:2,-0.0010);
    \node[anchor=north,inner sep=0pt,font=\normalsize,yshift=-2pt] at (axis cs:2,-0.001000) {60};
    \draw[line width=0.45pt] (axis cs:3,0) -- (axis cs:3,-0.0010);
    \node[anchor=north,inner sep=0pt,font=\normalsize,yshift=-2pt] at (axis cs:3,-0.001000) {70};
    \draw[line width=0.45pt] (axis cs:4,0) -- (axis cs:4,-0.0010);
    \node[anchor=north,inner sep=0pt,font=\normalsize,yshift=-2pt] at (axis cs:4,-0.001000) {80};
    \draw[line width=0.45pt] (axis cs:5,0) -- (axis cs:5,-0.0010);
    \node[anchor=north,inner sep=0pt,font=\normalsize,yshift=-2pt] at (axis cs:5,-0.001000) {90};
    \draw[line width=0.45pt] (axis cs:6,0) -- (axis cs:6,-0.0010);
    \node[anchor=north,inner sep=0pt,font=\normalsize,yshift=-2pt] at (axis cs:6,-0.001000) {100};
    \draw[line width=0.45pt] (axis cs:0,0.005) -- (axis cs:-0.3,0.005);
    \node[anchor=east,inner sep=0pt,font=\normalsize,xshift=-4pt] at (axis cs:-0.4,0.005) {$a$};
    \draw[line width=0.45pt] (axis cs:0,0.02) -- (axis cs:-0.3,0.02);
    \node[anchor=east,inner sep=0pt,font=\normalsize,xshift=-4pt] at (axis cs:-0.300000,0.02) {0.02};
    \draw[line width=0.45pt] (axis cs:0,0.03) -- (axis cs:-0.3,0.03);
    \node[anchor=east,inner sep=0pt,font=\normalsize,xshift=-4pt] at (axis cs:-0.300000,0.03) {0.03};
    \draw[line width=0.45pt] (axis cs:0,0.04) -- (axis cs:-0.3,0.04);
    \node[anchor=east,inner sep=0pt,font=\normalsize,xshift=-4pt] at (axis cs:-0.300000,0.04) {0.04};
\end{axis}
\end{tikzpicture}
